
\subsection{DAOS introduction}
Lofstead et al. present DAOS as a foundational redesign of storage for exascale systems, replacing POSIX file and directory abstractions with a distributed, object-based model built around containers, versioned objects, and epoch-based consistency mechanisms. The paper motivates DAOS as a response to scalability and metadata bottlenecks in traditional parallel file systems, emphasizing high concurrency, copy-on-write versioning, and efficient support for multi-dimensional scientific data. This work provides the architectural context and design rationale underpinning modern DAOS deployments, forming the basis for evaluating DAOS performance and scalability on leadership-class systems such as Aurora \cite{daos-exascale-storage2016sc}. 

DAOS introduction \cite{daos-intro2020}. Using DRAM as the first-layer storage instead of PMem \cite{daos-dram-intro2023}. DAOS as HPC storage \cite{daos-hpc-storage2023cluster}. DAOS repo \cite{daos-repo}. ARM64 DAOS client \cite{daos-arm2023SCAsia}.

\subsection{DAOS applications}
Using DAOS to speed up weather workflows \cite{daos2024weather_flow}. HDF5-DAOS \cite{hdf5-daos2022tpds}. Adios-DAOS \cite{daos-adios2024ISC}. DAOS as a burst buffer for EJFAT data streaming \cite{mei2024sc-poster}.

\subsection{Aurora introduction}
Aurora introduction \cite{aurora-intro2025}. Slingshot RPC enhancement \cite{aurora-daos-slingshot-rpc2025cug}.

\subsection{DAOS performance analysis}
{
\color{blue}
(Needs some comparison between our results to their results.)
}
IOR benchmarking with server and client sides scaling \cite{daos-scalability2023SCAsia}.

The 20-node Aurora DAOS testbed results \cite{DAOS-aurora-20-node-test2024SC-W}. 

DAOS performance exploring \cite{daos-perf-explore2024SC-W}.

DAOS vs Lustre \cite{daos-vs-lustre2022pdsw}.
